\section{Introduction}
    \subsection{Motivation and Context}

The progressive aging of the global population has led to a significant increase in social isolation among senior citizens. To address this, AI-driven conversational agents have emerged as a scalable tool for providing companionship and monitoring. However, the effectiveness of these interactions depends on the user's willingness to reveal personal information—a process known as self-disclosure.
 \vspace{2mm} \newline
In the field of psichology, understanding how and when an elderly user shares sensitive details (health, memories, or emotions) is critical. This project is motivated by the need to transform unstructured dialogue data from over 16,000 sessions into actionable insights. By applying Social Penetration Theory and Large Language Models (LLMs), we aim to establish a robust methodology to measure this trust-building process. From a telecommunications engineering perspective, this work addresses the challenge of extracting high-level semantic features to detect vulnerability and improve the quality of human-machine interaction in sensitive contexts.


    \subsection{Problem Statement: Automated Self-Disclosure Detection}

Despite the proliferation of conversational agents for elderly care, there is a lack of robust, automated methods to quantify the emotional and personal depth of these interactions. While current Natural Language Processing (NLP) techniques can handle intent recognition or sentiment analysis, they often fail to capture the nuances of self-disclosure.
\vspace{2mm} \newline
The core challenge lies in the unstructured nature of dialogue data and the subjective complexity of what constitutes a "sensitive" or "deep" revelation. Existing approaches often rely on rigid keyword matching or general-purpose sentiment scores, which are insufficient for identifying the thematic layers (health, relationships, personal history) defined by social psychology. Furthermore, the use of Large Language Models (LLMs) as evaluators introduces a new problem of methodological reliability: there is no consensus on whether binary classification, multi-level rubrics, or subtopic aggregation yields the most accurate representation of user trust. Therefore, this research addresses the need for a validated, theory-driven pipeline to transform raw conversational data into a structured metric of self-disclosure in the context of elderly-AI interaction.






    
    \subsection{Research Objectives}

The primary objective of this thesis is to design and evaluate a methodological pipeline for the automated detection and classification of self-disclosure in dialogues between elderly users and conversational AI, leveraging Large Language Models (LLMs) and unsupervised learning techniques.

To achieve this overarching goal, the following specific objectives have been established:

\begin{description}
    \item[O1: Feature Extraction and Scoring.] To implement and compare different scoring strategies— specifically binary versus rubric-based (0-5) scales—using a model to extract high-level semantic features from 16,000 conversational sessions.
    \item[O2: Methodological Comparison of Classification.] To evaluate the effectiveness of direct category classification versus a bottom-up approach that utilizes subtopic aggregation through the application of \textit{Hadamard products}.
    \item[O3: Unsupervised Pattern Identification.] To perform clustering (\textit{K-Means}) on the extracted features to identify thematic domains based on Social Penetration Theory, justifying the implementation of an $N+1$ cluster architecture.
    \item[O4: Empirical Validation.] To validate the model's performance against a human-annotated Gold Standard of 120 sessions, measuring inter-rater reliability via \textit{Krippendorff’s Alpha} and assessing classification accuracy through \textit{F1-score} and \textit{Matthews Correlation Coefficient (MCC)}.
\end{description}
        
    \subsection{Thesis Structure}


    This document is structured into seven chapters that guide the reader through the theoretical, technical, and evaluative phases of the research:

\begin{itemize}
    \item \textbf{Chapter 2: Theoretical Framework and State of the Art} explores the psychological foundations of Social Penetration Theory and the technical landscape of Large Language Models (LLMs) and unsupervised learning.
    
    \item \textbf{Chapter 3: Dataset Description and Preprocessing} details the characteristics of the conversational dataset and the protocols followed to ensure data privacy and RGPD compliance.
    
    \item \textbf{Chapter 4: Methodology} constitutes the technical core of the thesis, describing the experimental design for feature extraction, the mathematical application of the \textit{Hadamard product} for subtopic aggregation, and the $K$-Means clustering configuration.
    
    \item \textbf{Chapter 5: Analysis of Results} evaluates the outcome of the clustering process, providing a thematic characterization of the identified self-disclosure domains.
    
    \item \textbf{Chapter 6: Validation and Evaluation} assesses the model's performance by comparing it against an expert-annotated Gold Standard, using metrics such as Krippendorff’s Alpha.
    
    \item \textbf{Chapter 7: Conclusions and Future Work} synthesizes the findings, discusses the limitations of the study, and outlines potential applications in the field of elderly care.
\end{itemize}


\subsection{Social Penetration Theory}


Social Penetration Theory (SPT) provides a foundational framework for understanding how individuals reveal personal information to others, a process termed \textit{self-disclosure}. Often illustrated through the "onion metaphor," the theory suggests that personality is structured in layers, and relational development involves the systematic "peeling" of these layers through communication.

According to the theory, the disclosure process is defined by two primary dimensions:
\begin{enumerate}
    \item \textbf{Breadth:} The variety of topics covered during an interaction.
    \item \textbf{Depth:} The level of intimacy or sensitivity of the information shared within those topics.
\end{enumerate}

In the digital era, SPT has been widely adapted to study Human-Computer Interaction (HCI). In this thesis, the theory is utilized to categorize conversational data into distinct thematic domains. By mapping user utterances to specific layers—ranging from superficial biographical data to more intimate disclosures regarding \textit{Health}, \textit{Relationships}, and \textit{Personal Development}—we can quantify the level of trust and openness established between the elderly user and the conversational agent. This theoretical grounding justifies the selection of our classification categories and provides the necessary psychological context for the subsequent unsupervised learning phase.


\subsubsection{The Layered Structure of Disclosure}

Taylor and Altman describe the structure of personality as a series of concentric layers, ranging from the most visible to the most private aspects of an individual. The transition through these layers is what defines the "penetration" process:

\begin{itemize}
    \item \textbf{Peripheral Layer:} Contains "public" or observable data. It includes biographical facts (age, origin) and surface-level preferences. In elderly-agent interactions, this corresponds to the initial exchange of greetings and general information.
    \item \textbf{Intermediate Layer:} Encompasses personal opinions, social attitudes, and historical narratives. While still relatively safe to share, these topics reflect the individual's values and professional or social trajectory.
    \item \textbf{Core Layer:} Represents the most intimate and sensitive domain of the self. It includes deep-seated emotions, vulnerabilities, chronic health concerns, and significant life-changing events (such as widowhood or trauma). 
\end{itemize}

The systematic movement from the peripheral to the core layer is governed by the \textit{norm of reciprocity}, although in Human-Computer Interaction (HCI), this dynamic shifts as the user perceives the agent as a non-judgmental confidant. This study specifically targets the identification of shifts from the intermediate to the core layers, as they are the most indicative of the chatbot's success in fostering a supportive environment for the elderly.



The situational context and the nature of the communication channel significantly influence how self-disclosure is enacted. According to \cite{Bazarova2014}, self-disclosure in digital environments is not merely a direct reflection of relational depth but is also driven by specific functional motivations that vary across different platforms. Their research highlights that the technical affordances of a medium—such as perceived privacy and audience reach—shape the intimacy and characteristics of the shared information. 

In the specific context of this study, this functional approach is crucial: the interaction between an elderly user and a chatbot represents a unique, private communication channel. This perceived intimacy may lower the "social cost" of disclosure, encouraging users to reveal information from the \textit{core layer}—such as chronic health issues or deep personal history—more freely than they would in other digital social settings.


\newpage